Agradeço, primeiramente, aos meus pais, Roberto e Neila. Afinal, eles foram os responsáveis por toda a minha formação. Cada pedra em que pisei e que me trouxe até este momento teve uma mãozinha deles. Obrigado por nunca deixarem faltar nada à nossa família e por sempre me permitirem escolher o meu próprio caminho. Literalmente, sem vocês, eu não estaria aqui. Agradeço também à minha irmã, Camila, por ter sido a pioneira da família a estudar computação. Você certamente me influenciou na escolha dessa área, mesmo que talvez nem saiba disso. Também não poderia deixar de agradecer a todos os meus familiares que apoiaram os meus estudos e se preocuparam com o meu futuro. Eles ficaram tão felizes quanto eu quando ingressei na UFC e tenho certeza de que também celebram esta conquista. Muito obrigado a todos os meus familiares.

Agradeço à minha namorada e futura esposa, Marília, pelo apoio durante todo este processo, por demonstrar interesse pelo meu trabalho, por me escutar falar sobre ele em tantos momentos e por fazer perguntas que certamente contribuíram para fortalecê-lo. Seu apoio foi uma das grandes forças que me tiraram da inércia e me ajudaram a concluir este trabalho. Agradeço também à minha sogra, Cilene, companheira de pós-graduação. Foi muito divertido compartilhar esse momento com você. Juntos, vivenciamos dores e alegrias semelhantes durante este processo.

Agradeço ao meu orientador, Yuri, por aceitar me orientar mesmo quando eu ainda não possuía experiência prévia em pesquisa e por me conduzir a uma área de investigação tão interessante. Se não fosse por você, eu não estaria aqui para contar esta história. Agradeço ao meu grupo de pesquisa, CRAb, por disponibilizar as máquinas utilizadas na realização dos experimentos. Agradeço ao meu coorientador, Bento, e ao professor Creto por contribuírem continuamente para o aprimoramento da pesquisa. Agradeço ao Rodrigo, companheiro de pesquisa, por todas as nossas discussões nas quartas-feiras, por meio das quais conseguimos refinar cada vez mais o nosso trabalho. Agradeço ao Anderson por ter tido, no passado, a brilhante ideia do glaucoma. Agradeço também ao professor Gilvan por aceitar participar da banca examinadora deste trabalho. Suas contribuições foram muito importantes. Agradeço ao lendário professor Tarcísio por sua participação neste trabalho, tanto diretamente, como integrante da banca examinadora, quanto indiretamente, por ter inspirado tantas pessoas a ingressarem nessa área.

Agradeço aos servidores da UFC, especialmente a Vlademiro e Eldo, pela ajuda no laboratório, tanto com os equipamentos quanto com o acesso aos espaços. Agradeço também à equipe de limpeza, que sempre manteve não apenas o laboratório, mas todo o bloco limpo. Em especial, agradeço a Lourdes, por abrir a porta do LIA quando eu chegava cedo demais.

Agradeço às instituições de fomento à pesquisa pelo apoio durante o período do mestrado. O presente trabalho foi realizado com apoio da Coordenação de Aperfeiçoamento de Pessoal de Nível Superior -- Brasil (CAPES) -- Código de Financiamento 001. Por fim, agradeço ao LSBD pelo apoio no último semestre do mestrado, por meio da concessão de uma bolsa.